% chapters/06/evaluasi.tex
\bagian{Asesmen dan Instrumen Penilaian}
\label{sec:asesmen-fisika}

\begin{asesmenbox}
\textbf{Kedudukan asesmen bab ini.} Bab ini menuntaskan \SUBCPMK{1.3},
sehingga ketiga indikatornya diperiksa melalui penilaian artefak program
(Tugas Praktikum 04) dan soal konsep.

\vspace{0.4em}
Materi Bab~VI termasuk bahan \textbf{Kuis 1} pada minggu kelima (Materi 1--4),
yang berarti bab ini adalah materi terakhir yang diujikan pada kuis tersebut,
dan termasuk bahan \textbf{Ujian Tengah Semester} pada minggu kesepuluh.
\end{asesmenbox}

\subsection*{Instrumen 1: Rubrik Praktikum 04}

\begin{table}[htbp]
\centering
\caption{Rubrik penilaian Praktikum 04 beserta indikator yang diukur}
\label{tab:rubrik-praktikum-04}
{\setlength{\tabcolsep}{4pt}\footnotesize%
\begin{tabularx}{\textwidth}{@{}L{3.2cm} C{1.3cm} C{2.1cm} Y@{}}
\toprule
\headrow \textbf{Kriteria} & \textbf{Bobot} & \textbf{Indikator} & \textbf{Deskripsi capaian} \\
\midrule
Integrasi Euler dan lantai & 35\% & \ind{a} &
Kecepatan diperbarui sebelum posisi, benda berhenti di lantai tanpa terbenam
maupun menembus, dan koefisien restitusi pemain bernilai nol. \\
\addlinespace
Lompatan & 35\% & \ind{b} &
Lompatan hanya bekerja di tanah, menahan tombol tidak menghasilkan lompatan
berulang, dan tinggi lompatan terukur sesuai dengan hitungan. \\
\addlinespace
Pantulan & 30\% & \ind{c} &
Bola memantul di lantai dan di kedua dinding dengan tinggi yang menyusut,
pantulan kecil dihentikan, dan hasil percobaan koefisien restitusi tercatat
pada lembar kerja. \\
\bottomrule
\end{tabularx}}
\end{table}

\subsection*{Instrumen 2: Soal Konsep}

Jawablah dengan singkat. Kunci jawaban beserta indikator yang diukur
disajikan pada Tabel~\ref{tab:kunci-konsep-ch06}.

\begin{enumerate}
  \item Tuliskan dua baris integrasi Euler dan jelaskan urutan pengerjaannya.
  \item Mengapa benda dapat menembus lantai apabila pemeriksaannya memakai
        tanda kesamaan?
  \item Sebuah lompatan hendak dibuat setinggi 180 piksel dengan gravitasi
        2000 piksel per detik kuadrat. Hitunglah kecepatan lompatan yang
        diperlukan.
  \item Mengapa lompatan harus dibaca sebagai peristiwa, bukan sebagai keadaan
        tombol?
  \item Sebuah bola dijatuhkan dari ketinggian 120 piksel dengan koefisien
        restitusi 0,6. Hitunglah tinggi pantulan pertamanya.
\end{enumerate}

\begin{table}[htbp]
\centering
\caption{Kunci soal konsep Bab VI beserta indikator yang diukur}
\label{tab:kunci-konsep-ch06}
{\setlength{\tabcolsep}{4pt}\footnotesize%
\begin{tabularx}{\textwidth}{@{}C{0.8cm} C{2.1cm} Y@{}}
\toprule
\headrow \textbf{No} & \textbf{Indikator} & \textbf{Pokok jawaban} \\
\midrule
1 & \ind{a} &
\code{velocity.y += gravity * delta_seconds} lalu
\code{position = position + velocity * delta_seconds}. Kecepatan diperbarui
lebih dahulu agar posisi memakai kecepatan yang baru; apabila dibalik, benda
tampak jatuh lebih cepat daripada seharusnya. \\
\addlinespace
2 & \ind{a} &
Karena benda yang bergerak cepat melampaui lantai dalam satu bingkai, sehingga
posisinya tidak pernah tepat sama dengan nilai lantai. Pemeriksaan harus
memakai tanda ketidaksamaan, lalu posisinya dikembalikan ke permukaan lantai. \\
\addlinespace
3 & \ind{b} &
$v_0 = \sqrt{2gh} = \sqrt{2 \times 2000 \times 180} = \sqrt{720000} \approx 849$
piksel per detik. \\
\addlinespace
4 & \ind{b} &
Karena peristiwa \code{KEYDOWN} hanya terjadi sekali pada setiap penekanan,
sedangkan keadaan tombol bernilai benar selama tombol ditahan sehingga benda
akan melompat berulang-ulang dan seolah-olah terbang. \\
\addlinespace
5 & \ind{c} &
$h_1 = e^2 h = 0{,}6^2 \times 120 = 0{,}36 \times 120 = 43{,}2$ piksel. \\
\bottomrule
\end{tabularx}}
\end{table}
